\textbf{Turning Toward Refuge}

\textbf{Dhamma Reflections by Ajahn Karuṇadhammo}

\textbf{Contents}

Foreword

Making the Worst of Circumstances the Best of Practices

Know, Watch, Let Go

Me and Mine

Our World Through Perception

A Sublime Abiding

The Delusion of Forever

Turning Toward a True Refuge

\textbf{Foreword}

Ajahn Karuṇadhammo has been part of the Abhayagiri community since its founding in 1996. Before joining the community, he worked as a nurse in Seattle, during which time he encountered the Theravadin tradition. Ajahn Karuṇadhammo received full bhikkhu ordination in 1998 under Luang Por Pasanno. Throughout his monastic life, he has served the community in many ways, including as co-abbot.

Over the years, Ajahn Karuṇadhammo has generously shared Dhamma reflections with the many people who have come to practice and listen to his talks. Offered with warmth and clarity, his teachings have encouraged many to pause, reflect, and gently turn the heart toward what is true.~

The reflections gathered in this book are a selection based on Ajahn Karuṇadhammo's original talks, carefully edited for the written page while preserving their spirit and intention. These talks deeply resonated with us. \hl{As with other Dhamma teachings}, they continue to reveal new layers of meaning through mindful reflection and practice.~

For some time, it has been our aspiration to make these reflections available in written form. Most of Ajahn Karuṇadhammo's talks can be found online in audio form and continue to benefit listeners around the world. Reading, however, offers a different kind of engagement. It allows us to move at our own pace, to pause over a single sentence, to return to a passage that speaks to our present circumstances, or to revisit a reflection months or years later with fresh understanding.

Rather than presenting the Dhamma as a set of doctrines to be accepted, Ajahn Karuṇadhammo reminds us that ordinary experiences---moments of joy, disappointment, fear, loss, or confusion---can become the very ground for practice and opportunities for wisdom.~

His reflections encourage us to look more carefully at the experience already unfolding in front of us. They remind us that liberation is cultivated not somewhere else, but in the midst of the very conditions of our daily lives. Rather than offering simple answers, these teachings often leave us with questions to carry into our own practice, inviting continued investigation and deeper understanding.

This book itself is also the fruit of Dhamma in action. During the past year, many conditions came together to make its publication possible. The project grew from Suwalee Sea Kochawat's aspiration to bring these reflections together in book form. \hl{I} was fortunate to pick up the initial seeds from Sea and help coordinate the project through the transcribing and editing process. Jay Jacobs and Kristin Ohlson generously undertook the initial editing of the transcripts. Khantiko Bhikkhu offered further editing and proofreading, Hisayo Suzuki carefully attended to punctuation, copy editing, and proofreading, and Tyler Dennis provided formatting support. We are also deeply grateful to everyone who offered financial support to make the printing of this book possible.

It has been deeply inspiring to witness how so many individuals, each contributing in their own way, helped bring this project to completion. This collaborative effort is a beautiful expression of dāna, generosity, and spiritual friendship.

We offer our heartfelt gratitude to Ajahn Karuṇadhammo for his years of patient teaching, his guidance, and his example, and to everyone whose generosity has made this book possible. May these reflections offer support and encouragement to all who open these pages, and may they help us return, again and again, to the direct knowing of the Dhamma, where true peace and freedom are found.

With gratitude,

Anat Porat

Suwalee Sea Kochawat

\textbf{Making the Worst of Circumstances the Best of Practices}

Based on a talk given on November 4, 2021

We're in the middle of our Kathina celebration here at Abhayagiri, and Luang Por Pasanno and I have just returned from spending the Vassa at Pacific Hermitage. It\textquotesingle s really nice to be back here and to enjoy the camaraderie of bhikkhus, lay friends, supporters, and the broader community. We can rejoice in the flourishing of the good thing that's been happening here, including the many Dhamma reflections that have been offered.

At today's Patimokkha reflection, Luang Por Pasanno brought up the theme of rejoicing in this communal gathering of goodness, the goodness that comes with the coming together of people who are all committed to very high values. He noted that there's so much negativity and cynicism these days in the world around us. We see it in political and social issues. There are many difficulties fueled by the pandemic in the world these days, and people are feeling stressed. Their ways of responding feel natural, but are often marked by negativity and cynicism, which are not useful.~

We reflect on the importance of not getting caught in those states of mind. It's so easy to do. It's kind of a tunnel that's really easy to go down. We hear things that excite the mind, and we look for external situations on which to blame the difficulties that we experience. We think that if only we can right that wrong then everything would be okay. It's just a relentless, downhill mind state that doesn\textquotesingle t go anywhere if we buy into it.~

It's this way with other difficult, obstructive states of mind, too. Sometimes they are unrelenting, and even if we are trying to live a good and useful life, we can still feel bombarded. They can be very hard to bear. Whether it's anger, relentless sensual craving, doubts, confusion, delusion, or restlessness, to name just a few, it's so easy to let them expand into an even more negative or cynical mind state. If not projected outward to blame some external source, then they may be projected inwards with self-blaming, a sense of worthlessness, or an inability to make progress on the path.

Sometimes the difficulties, the obstructions that we experience in our practice, just come so hard at us. We get this sense of despair and we wonder, ``What's the point? Why do I put up with this? It's not going to go anywhere.''

We can experience great periods of doubt and difficulty. We can acknowledge that it's a common experience for all of us. Every one of us who has gone through periods like this, in one way or another. Knowing this doesn't necessarily make it any easier, but sometimes it's useful to expand beyond the self-referencing and realize that it's not just ``me.''

The question is, of course: How to deal with this? How to use the resources that we have to lift ourselves up out of it, to not stay down in a state of difficulty or despair? What resources can we muster to pull ourselves up and keep going?~

I've been reading a very compelling book that somebody gave to me recently. I'd heard of it before but had never taken the time to look into it and read it. It's called \emph{Man's Search for Meaning} by Viktor Frankl, and I've just finished the first section.

Viktor Frankl was a psychiatrist in Austria at the time of World War II. He was among those taken and interned into several concentration camps during the war and was a prisoner there for several years. The book contains his recollections about his state of mind while he was in the concentration camps. Right from the beginning, he was very clear that he wasn't writing this book as an account of all of the atrocities and painful circumstances that occurred while he was there, although, of course, he had to relate some of them to provide the context for his own internal processes. He then expounded on these states of mind and also what he saw as the internal mechanisms at work for the other people that he was interned with there.

He was an incredibly compelling and reflective person, and one of the few who was able to transcend the horrors of being in that situation. There were the physical horrors, the extreme deprivation, starvation, torture, extreme cold, hard work, malnutrition, barely enough to keep alive, and all of the diseases that would go around in these extremely crowded conditions.

This is if you were lucky enough to have escaped the extermination chambers, having watched your friends and relatives being taken already. In addition to physical deprivation, there were also the psychological difficulties, the horrors of being psychologically tortured day after day. He was able to stay on top of it enough to honestly experience and allow in what was happening, without moving into a state of denial or repression. He was also determined not to give in and let it break him. He noted a common pattern with many of the prisoners: after the initial horror and fears of being brought into the camps, there is a middle phase of intense indifference and apathy. This state of psychological annihilation, tuning out, seemed to be a way of escaping the emotional pain, just to be able to cope. He could see this happening around him and could see the allure of this coping mechanism, but he did not allow himself to completely collapse into that state of mind.

He had an inner resource, an inner drive to not give in to difficult, horrifying mind states. Over time, he was also able to be of service to his fellow inmates and would offer periods of reflection and support in ways that were under the radar to avoid getting punished and beaten.

It was in these subterranean ways that he was able to help others and to pull himself up time after time when doubts, confusion, and difficulties arose. It's a very inspiring story that we can all benefit from, reflecting on how he was able to transcend such an extreme situation.

There is this potential to not succumb, to not fall down deeply. Several of the ways that Frankl would do that would be to just not allow thought patterns dwelling on hopelessness and despair, to encourage memories of good things, to allow himself the hope of a future return to freedom and goodness. He would think of a return to more normal states of mind. He would allow his mind to attend to memories that were beneficial and would bring up a sense of hope and warmth, no matter how impossible they seemed. He would direct his thoughts in that direction rather than just succumbing to the present moment horrors of what was going on, without denying or repressing them, but also not dwelling on them.

Even though he didn't write about it in this way, one of his observations or realizations seemed very Buddhist to me. He described the patterns of stimulus and response that we all operate under and that the key to freedom and happiness is noticing and realizing the moment between stimulus and response. Whether it's coming from physical circumstances, mental, or psychological experiences, there is a gap between stimulus and response. There's a pause, no matter how small, and then the response.

Viktor Frankl's conclusion from his experience was that true freedom is found in being able to navigate that gap, so that we can choose where to go in any particular experience or situation. To me, that insight is the crux of our training, too. Before we respond, we can pause to acknowledge our experience, and then we can choose how we want to direct our attention. We do have that possibility.~

If he can do that in such an extreme situation, I have hope that all of us, myself included, can also do that when obstructions and difficulties come into our lives.

We should take delight in what we've been experiencing the past few days and in general how fortunate we are to have access to this kind of environment that supports the skillful attention towards goodness. Sometimes we can lose our perspective and forget that there are pockets of goodness---monasteries or our homes---where we can create an environment that supports reflection and turning towards the good. The environment supports turning towards the skillful and what is beneficial to ourselves and to other people. Places like this are repositories of skillful means and provide direction on where to turn, where to put our hearts and minds. They nurture the values that we can aspire to, the values that lead to an end of suffering. They help us develop the qualities of virtue, goodness, and generosity. We can cultivate the wisdom and compassion that is required in the difficult circumstances that we experience as human beings. Occasions like coming to the monastery for a period of time help us bring these resources back into our lives, back out into the world. This is the kind of world that we want to create. This is the kind of world that we can create, even in the midst of cynicism and negativity.

So really take refuge in the skillfulness and the good-heartedness of people who are struggling in the same ways that we are, trying to overcome these difficulties that are part of human experience. We all share this aspiration, and there's a set of skills and values that we can nourish and keep to the fore all of the time, with the proper effort. We don't have to succumb.

I've been reflecting on the importance of noticing when we are deeply involved with the perception of self and how this compounds and complicates our world and our experience. In a question-and-answer session with Ajahn Amaro, somebody asked him about the state of boredom. He responded, ``You can't be bored if you are dwelling with the perception of not-self.''~

When we are operating through and entangled with the perception of self, we can fall for obstructive mind states, whether it's boredom or one of the many others.

If we can learn to cultivate over time, little by little, the perception of not-self, we have the opportunity to see these obstructive habits and mind states as not who we are. Even though we might still experience them, we don't have to pick them up, own them, and compound them.

From a reflective position, we can understand: ``These are phenomena. These are selfless phenomena happening. It's painful. It's difficult. But I don't need to pick it up as mine.'' When we have that perspective, then we have energy and space to be with them. We have the energy and space to be in the gap between stimulus and response.

It takes the sting out of these obstructive mind states and quells the proliferative tendency. Dwelling in, refreshing, and reflecting on the perception of not-self is a major antidote to our experience of suffering. There are all the conventional ways of contemplating and developing the perception of not-self. It's also useful to contemplate the commonality of our experience with other people.~

Before I was a monk, I remember reading some accounts from Stephen Levine where he was recounting the experience of a woman who was in the process of dying. She was very ill with a terminal illness, and in addition to the relentless pain and physical discomfort, she was caught in mental resistance as well, with a lot of upset and anger. ``Why is this happening? Why can't I get over this? Why is this so painful?'' Due to the painful circumstances of her condition and the fact that it wasn't going to get any better, she was really suffering.

The struggle and the intense suffering was going on for quite some time. She was reacting to the bodily difficulties and struggling hard with her eventual demise. As this was happening, she was also trying to be reflective and conscious of the experience. At one point, it became so difficult that she went into a kind of transcendent state. She started to have visions and thoughts of all of the other people in the world who were going through similar experiences. She saw in great detail the struggling and suffering of person after person, with their own unique circumstances, time after time, in almost terrifying regularity. Through that experience she had an epiphany~of, ``This isn't me. This isn't mine. I\textquotesingle m not the only~one. This is not \emph{my} suffering. This is \emph{the} suffering.'' The personalization~of it as \emph{my} suffering fell away, and she realized~and saw the enormity of the situation. This is the condition~of the world. It wasn't hers only. It's a universal experience.

Being able to transcend the self-absorption, the self-obsession~that was a unique experience~happening to her was her liberating~moment, her epiphany. She saw it as \emph{the} suffering.~

I've always considered the universal experience~to be a really important thing~to reflect on for ourselves, too---not just the conventional~ways of trying to chip away~at this perception of self. We can realize with these experiences, these universal difficulties within the human realm: we're all in it, we're all doing it. It's the nature of the human condition. That opens up~a great sense of compassion, and we can let go~of the singularity of our own worlds.

Whether it's horrible experiences like being in a~concentration camp---the worst we can probably imagine in the human realm---or whether it's the internal struggles of our own~minds, we can realize that~this is the human condition and that none of us are exempt~from these experiences.~And it's~poignant. It's that very deep understanding and penetration of the First Noble Truth: there is suffering. With stories like Viktor Frankl's~and the stories of practitioners---nuns and monks, laywomen and laymen---going through these~experiences with honesty, clarity,~and openness, there comes the realization that~we don't have to succumb. We don't have to~buy into it. We can, with~effort and the support and encouragement of each other, realize that there's a way~beyond this.

There is an ending~to this if we want to see it through, if we attend~to it in the right way and nourish and~nurture the supportive~qualities in our life: the goodness that comes from~attention to~virtue, attention to generosity, attention to sharing~and helping each other. We share those~abilities as we share the suffering in our lives.~And it's this spiritual friendship that~is, as the Buddha said, the whole of the holy life, whether it's~the holy life in the monastery or the holy life in the householder's world. We have to maintain~a refuge in the Sangha, this support. We have to come together. We have to support~each other and be there for each other. We can't do it~without that support and that refuge.

I just feel so grateful to~be part of this community and to see the people coming day~after day. If ever I think, ``Am I doing a worthwhile thing? Is there~something else that I should be doing? Is there a better way to do this? Am I up to this?'' If I ever have those thoughts, I can reflect on the numbers of people that come here to support us and join in with our practice. People want us~to be here, all of us.

People want us~to keep doing what we're doing because we do help to keep that alternative alive in the world. We don't have to~succumb to the forces of difficulty, of~negativity, of cynicism. People appreciate being able to tap into that.

It takes a community of us to keep this alive. Laypeople, lay supporters, monastics, doing what we can do to keep this alive because~all the material things, all the ways of the world, those are all so fleeting. If we want to have a happy, peaceful~mind when it comes time for us to die and move on, then this is what will carry us through those periods.

This is the long-lasting~happiness that we can experience. It's done through Sangha, done through our~assembly of noble beings. As the chanting goes, ``They give occasion for incomparable goodness to arise in the world.'' Here we all are, keeping it alive. If we feel despair at times, then it's good to reflect~on what it is that we can do to help~keep this a sane place to live.~

\textbf{Know, Watch, Let Go}

Based on a talk given on March 28, 2019

This is the last Lunar Observance Day of our Winter Retreat, and we're about to~transition to moving back into more engagement---guests, visitors, work routines, the ``normal'' schedule of the monastery. We can watch the energy as it starts to build momentum towards that transition. Depending on the individual, there may have been a fair amount of settling, quietude, and gathering of energy during retreat, and once the restraints of noble silence and a reduced schedule are relaxed, the energy can start to bubble over.

I think we've all experienced that on retreats. When the noble silence is lifted and we're resuming normal activities, we find ourselves talking loud and fast, blurting out lots of runaway-train ideas. Watch this transition; see how it feels and be mindful.

There is a famous line from Ajahn Sucitto, just before a Winter Retreat started at Chithurst Monastery. Somebody asked him, probably a week or two before the retreat was supposed to ensue: ``Ajahn, when is the retreat going to start? What day is the retreat actually going to begin?''~His reply was: ``For some people, the retreat has already started. For some people, it will never begin.''~

We could probably say the same thing for the ending of the retreat. For some people, the retreat has already ended, or maybe never started. For others, it will continue on. That, in my mind, is the goal: to keep some of the retreat feeling. We can protect that feeling without clinging to it and without lamenting any change that does occur.~

I think most people here had a fruitful retreat, if not always easy. The community as a whole was committed to the practice and was definitely on retreat. We can work to maintain any skillful qualities that we've learned and to create habits to nurture and protect them. One way to do that is to notice the difference between outward-going momentum and the stillness~that has been nurtured through various ways in the past three months.~

During one of the readings a few days ago, there was a question about what the term ``\emph{saṅkhāra}'' refers to. Also, the word ``\emph{saṅkhāra}'' appears in the funeral chanting in the steps of dependent origination. It's the second factor: based on ignorance, \emph{saṅkhāras} arise. It also appears in the last reflection on impermanence that we do~as part of the funeral chanting: \emph{aniccā vata saṅkhārā}. It's a word that is used broadly, variously defined as conditions, fabrications, mind conditioners, volitional formations, and mental formations. It's used differently in different circumstances, but it has the quality of flow to it.

Whether it's used as a general reference to all conditioned phenomena,~or to more specific phenomena, it has the quality of flow, of movement, rolling on. It's very useful to view it as a fluid experience~rather than trying to pin it down as a static object. That can be difficult to do. We often speak in terms of nouns: \emph{Aniccā vata saṅkhārā}, ``all compounded things are impermanent.'' \emph{Saṅkhāras}---compounded things, creations, mind states---those nouns tend to solidify our perception of saṅkhāras.

All of these mental phenomena have the quality of movement and flow. Speaking from the perspective of experience, there is nothing that's solid. It's all an experience of movement and flow. The quality of \emph{cetanā}---intention, volition---is one of saṅkhāra's primary characteristics. Intention, the mental impulse or driving force behind all action, feeds the saṅkhāra process.~

Saṅkhāra, in its active phase, produces mental qualities and activity. It serves as the basis for action that manifests through body, speech, or mind.~

Sometimes saṅkhāras are experienced as \emph{vipāka}, resultant conditions. Perceptions and feelings are the end result of the intentional process. They are resultant conditions, but are still part of the saṅkhāra flow. This rolling on of productive and resultant experience conditions even further intentional action.~

If the process is driven by greed, hatred, or delusion, then we keep creating formations that support a trend or habit. This can create a mental movement towards or away from an object. There may be confusion and ignorance regarding the activities that we're producing~through body, speech, and mind. Those activities, saṅkhāras\emph{,} further manifest in ways~that reinforce old patterns of perception~and the accompanying feeling tone.

The feeling tone of unpleasant, pleasant, or neutral,~conditions the next uprising of intentional activity, creating more saṅkhāras\emph{.} The flow keeps going. This endless motion keeps propelling us through time and experience. It's rare that we allow things to slow down and stop.

Another word that we've been talking about recently~is \emph{papañca}. In a discussion with Anuruddha, the Buddha lists eight thoughts of a great being (AN 8.30). In this teaching, the Buddha says, ``This Dhamma is for one who delights in \emph{nippapañca}, not for one who delights in \emph{papañca}.''

There are various ways to translate the word ``\emph{papañca,}'' including objectification, proliferation, or complication. We engage in papañca with the sense of I, me, or mine: the sense of self being at the base of it.

Papañca can be seen as the emperor of \emph{vaci-saṅkhāra}, verbal formations or verbal conditioners. Papañca is like vaci-saṅkhāra~blown up on steroids and fueled by greed, hatred, and delusion. When there is a strong \emph{kilesa} (afflictive state) in the mind and it engages with the thinking process, it causes a runaway-train effect. It accelerates, compounds, and feeds on itself. It becomes its own beast that then turns back on us.

We become the victim of this beast that has been created~through clinging and proliferation. What starts out as a perception and feeling based on contact with the outside world is complicated by thinking and proliferation. The process is reinforced over and over again.~

With my background as a registered nurse, I see cancer as a metaphor to this process of proliferation. Cancer starts with one aberrant cell that gets fed and nurtured~and is not recognized and destroyed by the immune system. It's out of the range of detection. It starts to reproduce and grow and proliferate. In the same way, our minds do this when there's no mindfulness. When there's no oversight, mindfulness isn't present to pick up on the little errant thought~that arises out of aversion, craving, or~confusion. We feed it and it starts to increase and proliferate. It takes over, like when cancer starts locally and then metastasizes to other parts of the body. The flow of thoughts, feelings, and emotions suddenly goes on a tangent and~lands and grows on something else: another memory, injustice,~or form of craving.~The whole process is self-reinforcing. We've got a fulminating disease, as they would say in the medical world.

On the mental level, these kinds of uncontrollable processes are~saṅkhāras, papañca as the extreme manifestation. We are trying to get a handle on~and see them so we can learn not to be the helpless victim. The Noble Eightfold Path provides the means for curing the disease.

Through sustained practice, we can understand how to let that process slow down and come to stillness. In the treatment of cancer, sometimes attempts to annihilate~and control the cancer cells are effective. Sometimes the effect is only temporary because we haven't completely eliminated the cancer cells.~Although we've suppressed them for a while, they come back. Similarly, we can try to annihilate our outward-flowing papañca~and place strong holds on it. This may work for periods of time, but we're really trying to get down to the root causes. We use \emph{sati} \emph{sampajañña}, mindfulness and clear comprehension, to stimulate our mental immune system in order to recognize and turn away from the causes~that start the whole process.

Part of rooting out the causes is learning how to slow down---to recognize and let go of the underlying tendencies. We try to stop reacting to, feeding, enabling,~and encouraging them. It's a process of release, of recognizing and letting go. ``Knowing and letting go'' as Ajahn Chah would say. We let go not out of resistance or aversion or pushing away,~but because we see that it's painful to hold on~to something that's hot and burning.

If we put our hands on a hot burning object, we immediately release it. It is a reflex. Similarly, we learn to recognize the habits which burn us when we grasp onto them,~and we let them go. The more we practice, the quicker we let go. We recognize the true nature of formations, where they come from and how they pass away. Aniccā vata saṅkhārā. All compounded things are impermanent.

Thinking of saṅkhāra as a flow rather than an object gives a sense of movement. Saṅkhāra is something that's in motion and is inconstant. However, when there's an experience of pleasure, we want to grasp on to it. We want to pin it down, solidify it, and keep it. We do that even with experiences that are uncomfortable.~

Grasping is based on a perception of ``me,'' a self that is trying to hold on and protect. The experience gels, hardens, ``firms up,'' as Ajahn Sucitto says, into what feels like a solid reality or a solid perception. At the time, we want to hold on to the familiar, even if it's unpleasant, difficult, or harmful because we seek stability. But if we see saṅkhāra as constantly flowing and changing no matter what we try to do, we eventually realize the futility of trying to hold onto it or push it away. We relax into letting go. We allow what is always changing to do just that, according to its own nature.

This is something that we see more clearly with a still mind. If we don\textquotesingle t jump on the moving train and go with all the sights, sounds, smells, tastes, tactile sensations, and thoughts, then we are more with the natural flow of conditions. We stand beside the track and watch the train pass by.~

``\emph{Uppajjitvā nirujjhanti}'' is another line from the aforementioned reflection, and is translated in different ways. Ajahn Geoff (Ṭhānissaro Bhikkhu) translates it as, ``They (fabrications or compounded things) disband as they are arising.'' It points to the momentary nature of all experience. When the mind is quiet and steps back to observe the flow, arising and passing away are not timebound. Passing away is inherent with arising. Arising and passing away are part of the same experience and to see this deeply is liberating because the natural impulse is to let go.

If all conditions are constantly changing, there's no need to create all the formations that help support the illusion of solidity. When we can do that, in that stilling there is peace. The flow of conditions still goes on, but the mind steps back. There is just the knowing of the flow of experience, the basic function of citta: the quality of awareness.~

How can we protect and maintain some of that stillness as we turn back into daily activity? We start with the basics: keeping our precepts.

If we want to have a quieter mind, without the regrets and disturbances that arise from harmful or agitating actions, following moral precepts is our foundation. It's our front line of protection. We cherish the life of even the smallest of creatures, such as mosquitoes or garden pests. We cherish life enough so that although we may try to get them to move away or protect ourselves from them, we don't annihilate them. Reflecting on gentleness and caring for sentient beings creates ease in the mind, a sense of quiet and stillness through non-harming.~

It's the same with all the other precepts. The second precept is to not take that which is not given to us, to not steal anything. We don't have to suffer with the remorse of that action, even on small levels. For example, we might think nobody's going to notice some minor form of theft that ``everybody does,'' like cheating on taxes. However, we don't want to have that niggling in the back of the mind.

The third precept is refraining from sexual misconduct, not taking advantage of others in a sensual way. The fourth precept is refraining from lying. Breaking these precepts will bring agitation and discomfort to the mind. Being careful to protect truthfulness and other wholesome qualities of speech is important to our peace of mind. We refrain from harsh speech, backbiting and even \emph{samphappalāpa}, nattering and purposeless chatter. When we follow these wholesome courses of action, we experience a greater level of settledness and stillness.

We can also find ways to simplify our lives. It takes more effort outside of the monastery because not everybody else is doing it. We can question how much we need and ask ourselves where the acquisitive mind takes us. The world is getting more and more complicated with all the various forms of stimulation and distraction that seem to be shortening our attention span, day by day. It's incredible. In waiting rooms, airports, and even walking down the street, people are totally absorbed in their electronics, listening to something with their earbuds. We can reflect on how this affects us and perhaps learn to not shut off from experiencing the world around us in self-absorption.~

Take time to ask yourself, ``What am I doing? Do I need to do this? How's this affecting my mind? How's this affecting the quality of my attention? Is it leading to inner stillness? Is it supporting mindfulness and clear comprehension?'' Take a break. Step back and remind yourself of the state of mind you want.

It's easy to get involved in work, activity, and stimulation. Doing things, picking up the pace and saying, ``Okay, in three months' time I'll get to go on retreat again.'' We go back and forth between living an intensely stimulating, distracting life and shutting it all off with special conditions, even coming to the monastery for a period of time. That's good to do. It's better than not doing it. But try to be consistent with the practice in daily life and really question the level of distraction you engage in.~

If you're in the car as a passenger, put down the phone for half an hour and just sit quietly. Try to do that. Take a break.

If we're intent on preserving the quality of stillness, work with sense restraint, protecting how we attend to the outside world. We do have to engage with and experience sights, sounds, and tactile sensations, but \emph{how} we engage is up to us. Without putting ourselves in sensory deprivation tanks and walling ourselves off from every possible kind of stimulation, we can watch how we attend so that we don't get fooled by the exaggerated qualities of the world around us. Whether it's the exaggerated attractions or aversions, we don't have to buy into the mask, the illusion, the enchantment of sights, sounds, and objects that catch our attention. Instead, we look at them with wisdom, in their true nature.~

We can absorb into objects that we find attractive because we attend to signs and features that stimulate and support desire and craving. We can counteract that absorption by looking at things more completely, looking at things in their context and getting a true picture of what we find so attractive. Is it really inherently attractive? What is it that seems to be attractive? What is desire going out to look for? Notice that process and give it a reality check with wise attention.~

Desire exaggerates the beautiful qualities of an object. It holds on to the object and builds on it. We can stop and say, ``Oh, I am creating a reality that's not really there.'' It's the same with aversion and dislike, how we create and compound stories, one on top of another, through inappropriate attention. We attend to the undesirable qualities and try to get away from them. We create a whole world around this perception. We can think, ``Maybe there are some good qualities of that person whom I find irritating or difficult.'' Try to see the whole person.~

Remember the human qualities we all have, including the undesirable or irritating qualities. We can realize, ``Okay, they're not circumscribed to one being who happens to be in my field of attention right now. These are qualities that arise and pass away, too.'' We don't have to hold on to them or run away from them. But we have to protect our minds from going out and getting lost in them.

We see the danger in lack of sense restraint. We step back and see the flow of the world, and we don't jump into it. We try not to get entangled. We can't stop the flow of \emph{saṃsāra}, but we can realize peace by dropping back and resting in the knowing mind, in awareness. Maybe we've had a few momentary glimpses of peace from our dhamma practice that give us a little bit of insight and inspiration: ``Maybe what the Buddha says is true. Maybe there is a path that can go to the complete cessation of entanglement.''~

Upāsikā Kee Nanayon described this kind of knowing. ``An inward-staying, unentangled knowing, all outward-going knowing cast aside.''\footnote{Upāsikā Kee Nanayon, \emph{An Unentangled Knowing: The Teachings of a Thai Buddhist Lay Woman}, trans. Thānissaro Bhikkhu (Valley Center, CA: Dhamma Dana Publications, 1995), 29, dhammatalks.org.} Our awareness moves inward: knowing, watching, letting go, and not jumping into the flow.~

In that complete stilling is true peace.~

\textbf{Me and Mine}

Based on a talk given on December 16, 2022

A few afternoons ago, I was having tea with some of the guests here at the monastery during the teatime Q\&A, and one of the topics that came up in our discussion was how to use certain teachings that we don't completely understand yet. How do we hold them when we're unsure about and haven't really experienced them ourselves or know what to think about them in relation to the teachings, the training? How do we use teachings as contemplative strategies rather than doctrines to believe in?

I thought that would be an interesting thing to explore a little bit tonight because it's a fundamental aspect of our practice in the way that the Buddha set it out. We aren't instructed to pick up teachings in a way that's just on a scholarly or intellectual level in order to present our views and opinions as being right and everybody else's view as being wrong. The Buddha pointed out the unskillfulness of that and how it doesn't actually get us anywhere other than ensnared and entangled with views, opinions, and conflict.~

So how do we use teachings for the purpose of ultimately realizing the end of suffering? Using them as contemplative strategies is generally how he instructs us to hold the practice, the teachings, and the discipline, both the monastic discipline and the lay precepts.

The word ``\emph{opanayiko}'' that we see in the chanting means ``to be brought in.'' The teachings are to be brought in, not just to be learned on an intellectual level. It doesn't mean that the Buddha doesn't have some categorical teachings that he presents as facts or as accurate representation of natural laws. For example, the teatime discussion that we were having was particularly around kamma and rebirth.

There are a number of teachings where he categorically says, ``As I understand it, as I've experienced it, this is how it works.'' The teachings in the category of Right View fall into that category: kamma and rebirth being two of them.~The Four Noble Truths is also a categorical teaching. But even with categorical teachings, they're not to be picked up just for the sake of understanding on a conceptual level but are to be applied as a strategy for reaching the end of suffering.

Both categorical teachings and other teachings that aren't categorical but are more reflective should be used as skillful means. As for the ones that we don\textquotesingle t understand or haven't experienced for ourselves, we can pick them up and use them as working hypotheses. We haven't necessarily proven them to ourselves, but we can pick them up on faith, not blind faith, but with the thought, ``Okay, what I've picked up and applied before from the Buddha's teachings seem to have had an effect. They seem to work. The Buddha also taught \emph{this}, which I haven't really experienced. I'll pick it up on faith because a lot of everything else that he taught has worked so far, and I'll try it out. I'll pick it up as a hypothesis, try it out, see what happens. What are the effects in my mind, in my heart?''

It's a learning process. It's an educational process that we're picking up. And we're using it with a spirit of inquiry, not in the spirit of doctrinal belief. One of the best examples of this that sometimes gets misunderstood, particularly when one first encounters the teachings, is what sometimes is referred to as the doctrine of not-self, \emph{anatta}, or the characteristic of not-self. Just the way it's phrased sometimes, for example using the word ``doctrine'' or ``characteristic,'' starts to turn it into something that is more along the lines of a description of a belief, rather than as a contemplative strategy. And I think it's really useful to undo that habit because a lot of people who are trying to understand this teaching end up with questions such as, ``Well, the Buddha says there is no self. If there is no self, what is all this: this body, this mind?'' And then from that comes a whole barrage of questions that just run in circles and don't lead to any kind of deep understanding.~

I think about a time quite a few years ago when I was staying at Pacific Hermitage in White Salmon, Washington, and Ajahn Sudanto and I were invited to give a presentation at a local community college on the other side of the river in Oregon. They were studying the world's religions and asked us to come and share some of our understanding of Buddhist practice. It was a great class, very lively people, with very intelligent, inquisitive minds. Eventually it came around to this particular teaching on not-self, and one young woman was asking a lot of questions about it. I explained how the teaching on not-self could be used as a contemplative strategy rather than as a doctrine of belief, as a way of deconstructing our sense of self.~

I was trying to point out that the sense of self is merely a perception, but that we can run into problems when we say, ``There is no self.'' If we start categorically saying things like ``there is no self,'' then it can lead to confusion because the sense of self is such a strong perception, which most of us experience. It really feels like there's me here. I'm here. I'm this body. I'm this mind. Or, I am in control of this body or in control of this mind. The Buddha didn't declare, ``There is no self.'' There's a line in our morning chanting, ``\emph{sabbe dhammā anattā,}'' and it's translated there as, ``There is no self in the created or the uncreated,'' which might lead to an annihilationist viewpoint. ``\emph{Sabbe dhammā anattā''} might be better translated as, ``All \emph{dhammas}, or all phenomena, are not self.'' We can see there's a difference between saying ``there is no self'' and ``all \emph{dhammas} are not self.'' The second translation refers more to that contemplative strategy.

During this discussion, I was explaining to her that as a contemplative strategy, we can go through what we normally think of as self. This is laid out in very succinct phrases in the \emph{Anattalakkhaṇa Sutta} as the five aggregates: form, feeling, perception, mental formations, and consciousness. In some form or combination of those aggregates, singularly or collectively, we impute a sense of self onto them. We take this body and/or this mind in its various forms to be who we are. The \emph{Anattalakkhaṇa Sutta} goes through each one of the aggregates and starts to unravel that perception with the Buddha's skillful questioning of the bhikkhus.~

I described how we go through each of these aggregates. We look at the body. We investigate it, and we see it very clearly through a pristine lens of \emph{samādhi}, a collected mind that's been settled in meditation. We start to examine our assumptions of this body being who we are or something that we possess. We carefully examine it piece by piece, part by part, function by function, and ask ourselves, ``Is this something that I have ultimate control over? Can I say that this is who I am or something that I possess?''~

We continue that investigation through the aspects of mind, including our feelings, the pleasantness or unpleasantness of experience, and the perceptions that we create. Perceptions are patterns that try to make sense of the world around us, based on data that comes in through all of the different sense bases: eye, ear, nose, tongue, body, and mind. We create mental perceptions to help give meaning to our experience.

We ask the questions: ``Are they who we really are? Are they completely under our control? Are they unchanging?'' And so on with thoughts, mental formations, and sensory consciousness.~

I was explaining this contemplation of going through them one by one with a clear mind and examining these aggregates to see whether they stand up to the test of something that we could call ``me'' or ``mine.'' The process is one of deconstructing our patterns of perception that impute a sense of self onto these five aggregates. It's a process of trying to uncover that this isn't who I am.

She was really intrigued by this and was processing out loud. She said, ``So, I think I get this. You go through this process of deconstructing who it is that you \emph{think} you are but really are not, as a contemplative strategy. And when you finally penetrate and realize that you really aren't those things, then what appears is who you \emph{really} are.''~

And I paused for a moment as I was thinking, and then my response was, ``Well, not exactly. You go through this process of deconstructing the sense of self in relation to forms, feelings, perceptions, mental formations, and consciousness. The result is finding out who you really \emph{aren't}. There is no second stage of finding out who you really are.''

It's just that one has completely discovered who one really is not. And then leaves it there to see what happens in the mind. This is using the teaching on not-self as a contemplative strategy.

As the teachings of the Buddha point out, with the complete understanding through contemplation, we develop a healthy sense of disenchantment with all of these things that we've formerly assumed to be who we are. We don't look for a sense of security, a sense of stability in these aggregates that are constantly changing. That's what then leads us to relinquishment, to letting go of constantly moving within those realms of trying to create happiness, stability, permanence, and security in the world of the aggregates. That offers the opportunity to open up to what the Buddha described as the Unconditioned, that which does not participate in or partake of the world of the constructed, the fabricated. However, he doesn't go into redefining it and reifying it or concretizing any kind of experience into a sense of self.~

There are many strategies similar to this particular contemplation that the Buddha gives to us; many teachings, many suttas that describe this process. There's a particular one that I've always liked. It's a very straightforward, simple simile where the Buddha is with a group of his disciples in Jeta's Grove, Anāthapiṇḍika's Park, a place where he spent a lot of time teaching.

He says, and I paraphrase, ``In this Jeta's Grove, if people were going around and collecting all the grass, leaves, twigs, and branches scattered around, and starting to walk off with them to take them to use them for something or to burn them, would you say, `Stop, stop, don't take any more of those twigs! Don't take us! Don't take me! Those twigs and grass and leaves, don't take me away! Don't burn me!' Would you say that?''

And the monks respond, ``Oh, of course not. And the Buddha says, ``Why is that?''~

``Because that doesn't have anything to do with us. That's not who we are. Those are just products of the forest.''

He said, ``Well, in the same way, you should not assume to take possession of, take hold of, to identify with that which isn't yours.

And what is it that isn't yours? Form, this body, isn't yours. Abandon it. Feeling isn't yours. Abandon it. Perceptions are not yours. Abandon them. Mental formations aren't yours. Abandon them. Consciousness isn\textquotesingle t yours. Abandon it. In the same way that you abandon and don\textquotesingle t take possession of and don't concern yourself with the forest products being taken away, don't regard any of these five aggregates as yourself. Abandon them. Abandon that which isn't yours.''

~

We can use this as a contemplative strategy to think, ``In the same way that all these bits and pieces of the forest are just mere elements, earth element in particular, we can regard the body as just mere elements: the organs, the fluids, the winds, the different aspects of heat, the four physical elements of this body. That's all it is. It has nothing to do with me or mine.''

There are many other suttas that work along similar lines to help deconstruct the sense of self. Pick them up, not as a doctrine, but as something to investigate.~

Another teaching that I found really useful is something that Ajahn Sucitto has talked about in reference to the body and the mind. He talks about them in terms of intelligence. Normally, when we think of intelligence, we think of it in regard to the thinking mind that is involved with conceptualization, abstraction, and planning---very active mental activities. Often when we think of an intelligent person, he or she is someone who's able to process information and make associations well. It's a very useful tool that we have, this intellect and the intelligence within it. It helps us to organize data and store data in our memories so that if it's useful data, we can then retrieve it and help to organize our world in useful ways.

If it's not used skillfully, it can also be co-opted very easily by the defilements of desire, aversion, and delusion. It can be put to harmful ends, creating all sorts of stories and unskillful states of mind, and then we dwell and proliferate on those. Whether that intellectual ability is used to our benefit or not depends on what's in charge at any one time. It can be a very powerful tool.

There are also other kinds of intelligence that we don't really consider very much, such as bodily intelligence.~ We can observe the body as a very independent organism that has its own innate intelligence, which we don't pay much attention to. We generally go on automatic pilot with our body. Say you're sitting down reading a book, and all of a sudden you get hungry. You put the book down, get up, and go to the fridge. You pull out a bagel. You cut it in half without cutting your finger off, pop it into the toaster, turn around, go find some mayonnaise or mustard, slather it on the bagel, add some cheese, some sprouts, put it all together, take it back to the table, and sit down.

You're probably not very conscious of the activities that you're doing. You may be talking to a friend. You\textquotesingle ve got your earbuds in. You're talking to somebody over the phone, listening to some music, or listening to a Dhamma talk on not-self while you're making your sandwich. You've gone through this whole process, and you haven't had to think about it consciously. You haven\textquotesingle t had to instruct your body how to get up, lift its foot, and move it forward, and then do the same with the other foot to go in the direction you want; how to open the fridge; how to cut the bagel; how to construct a sandwich. You're probably doing this without much awareness at all. The body is just doing all these activities that require incredible dexterity and precision. Meanwhile, you're talking to your friend and not even paying any attention to what you're doing. However, there's a general, vague awareness of motion, of something happening.

This is the bodily intelligence. The body has its own innate quality that helps it function, whether we're aware of it or not, whether we're conscious or not. It even functions when we're completely unconscious, like when under anesthesia during surgery; the body keeps on functioning. The heart keeps beating. The lungs keep breathing. The metabolic processes keep happening. When we are conscious, the body's movements and activities are done without much awareness. This is the innate bodily intelligence, and it's quite profound.~

This is similar with the heart/mind---not the intellect, but the \emph{citta}. \emph{Citta} is a process with its own innate intelligence that feels, that is sensitive. It intuitively picks up and experiences perceptions and feelings.

For example, something comes along that's a moving experience. We see somebody who is having difficulty, struggling or suffering in some way, and it doesn't take any kind of thinking process to feel a response. The citta can automatically feel compassion, a sense of caring and concern. It's an intelligence of the citta, of the heart. Rational prompting is unnecessary to understand that kind of experience.~

Or we do some sort of unskillful action that causes difficulty or harm, something we say or do that results in pain. If it's seen clearly and recognized, the heart responds automatically with a sense of, ``Oh God, why did I do that? That's so painful to see. I feel really bad about that.'' We don't have to go through a big intellectual process of analyzing step-by-step what happened. If we have cultivated a sense of conscience, this sense of skillful regret arises because the heart is impacted. This is the innate intelligence that our citta operates with, for better or for worse.~

There are these kinds of intelligence that if we really bring them into regular contemplation and experience them, we can start to realize, ``Oh wow, my thoughts, my feelings, my perceptions, even my bodily activities keep on operating, whether I am conscious of them or not, whether I am in control or not, whether I'm present or not.'' And we can start to step back a little bit from this quality of our mind enmeshed with the object, where we identify so strongly with the body and the activities of the mind, even the very subtle fluctuations of the mind. All of these are happening on their own steam, without the presence of an agent that is in control. When we can use our meditation to quiet the mind and then start to examine these experiences, not thinking about them so much, but watching them as they manifest, with a little bit of growing distance between the quality of pure knowing and the activity itself, then this can give us a sense of this perception of not-self. With this perception, there can then be a little bit of disentanglement, deconstruction from this marriage to a sense of ``I am.''

I have found that these are useful tools. I have a lot of appreciation for the people who have articulated the teachings in this way. It brings the practice out of the realm of doctrinal belief and into the realm of a strategy that we can use, and that results in deep insight such that our perspective of who we are in this world starts to change. We discover ways to help free the heart from the entanglement that we've woven for ourselves.

When we do that, as I mentioned before, there arises a sense of healthy disenchantment. We can clearly see the ridiculousness of how we keep creating all of these paradigms in a way that we think will give us some sort of long-term gratification but doesn't. It's short term. Instead, we adopt new strategies to help steer the mind towards that clear seeing, an awareness of how much we construct this perception of self to provide a sense of security for ourselves. We don't need to do that. As we are able to realize, understand, and experience that deeply, then a profound sense of freedom can arise, an appreciation of letting go and relinquishing these old patterns. This can serve as the stepping stone into what is referred to as the entry into emptiness, the experience of emptiness of self.

There is a freedom that results from that, and we are able to use that as our point of leaping off into an experience of the Unconstructed, the Non-Manifest. We can directly realize the Four Noble Truths: understanding suffering, understanding how it's caused, and practicing the way leading to the complete ending of suffering. That's the potential from using all of these teachings as contemplative strategies rather than just as interesting doctrines to believe in.~

\textbf{Our World Through Perception}

Based on a talk given on January 16, 2023

Yesterday was the anniversary of my mother's death six years ago, and I spent a little bit of time today looking at some photos. For the memorial service, my brother, sister, and I, along with various other family members, had put together a slideshow---a brief synopsis of her life in photos. I was struck once again at the variety of images, perceptions, and people that constitute what it is that I call my ``mother.''

There were photos of her as a young girl, of her growing up into young adulthood, and photos of her on the beach, probably in her twenties. She had dark brown hair and was a fairly attractive young woman. There were some wedding photos and some of my mother and father together in their early years. She started to have kids, and there were pictures of her feeding my brother as a baby. She was looking a little older and maybe a bit tired, but still lovingly feeding him. Then there were photos of her and the three of us, my brother, my sister, and I, as very young kids. She was getting into her thirties.

As she got older, into her middle and later years, the photos were of her with her grandchildren and my sister, and some of myself in my thirties. The photos then progressed through her aging, final years, illness, and finally a number of photos of her after she passed away in her bed in the nursing home, where she spent the last six weeks of her life.~

I found myself contemplating, ``Who was she? Who was this woman named Glenn DeMaria?'' (I was named after her, which caused a lot of confusion for people).

Who was this person that we label as ``Glenn DeMaria''? Was she that young girl on the horse, the mother with three children, the older woman in the wheelchair in her final days, or the lifeless corpse on a hospital bed? The worlds that we live in are worlds of perceptions, labels, and constructions of supposed realities. The people that inhabit these worlds are the same, in terms of experience. They are perceptions.

We create our perceptions of the world around us, and then we inhabit them. We take them all as so solid, so real. But right now, the entire lifetime of some person is just a memory.

We can think about that in relation to ourselves as well. ``Who am I? Who was I?'' Looking at some of those photos, there was a young kid, maybe three or four years old, with a completely different body than this one. Blonde, straight hair in the early years, curly hair and a bushy beard later on, and then a body that morphed into middle age. Then there were some photos of me with my mother in her last few weeks, when I was a sixty-one-year-old monk. I think, ``Okay, well, who was that? Where is he now? What was that all about?'' It's just a memory now.

You can ask that of yourself. Think back to your early life. Go through the various stages.You have a completely different body now. Body cells replace themselves more or less every seven years, so you don't even have the same body that you had seven years ago. Who was Lee Mintz? I remember a Lee Mintz who had big, bushy, dark hair and was fifteen or twenty years younger than he is now. Who was he before that? Or Ty Krug? Or Nick Teal? Josh Dultz? Who were you? Where were you? I don't even remember Tan Sampanno's former name. I just remember a Korean monk who showed up at Abhayagiri. Where did they go? Just memories. I can remember Debbie Stamp from at least thirty years ago. Where did she go? These are questions for reflection that challenge our assumptions of the world we create around us. How real is it? How solid is it? We take it on as so real, so solid, but everything before this very moment is just a memory. The future is just a projection. Experience is constantly moving, changing, shifting.

Why do we take it all so seriously? For my mother, there are no more memories. Assuming she has been reborn, ``she'' is experiencing another existence. Six years later, if she has been reborn into the human realm, she is sitting somewhere as a six-year-old child. She probably has no memory of Glenn DeMaria.~

Many Buddhist teachings instruct us to take a look at our perceptual world, the reality that we've assumed into existence. Those teachings encourage us to stop identifying with all of these constructs, to let the sense of ownership that comes with identification fade away. When we take on those constructs as who we are, that's the basis for suffering, \emph{dukkha}.~

The Buddha's teachings point towards recognizing the arising of dukkha. How can we reshape our perceptions such that we can relax the process of grasping that is based on mistaken perceptions? Perceptions are based on memory. Perception, the aggregate \emph{``saññā,''} is something with which we often strongly identify. ``Recognition'' might be a better translation in that all of our perceptions, the ways that we experience the world, are through filters. Filters cause us to remember certain related experiences from the past that we label and categorize to create a sense of continuity. We need perceptions to navigate through the world without going insane because of lots of random, sensory stimulation. We assign meaning to memories and experiences. ``Felt meaning'' is another translation of \emph{saññā} that I have heard. Recognition is a cognizing process. It's based on memories that help us to create categories of people, situations, and mental structures. That's what we do to get through life without losing our sanity. But the perceptions we create are up to us, and so many times we create perceptual patterns that result in suffering.

Some of the teachings that the Buddha gives us are approaches to reconstruct perceptions in ways that are more suitable for living a happy, contented life, for letting go of dukkha. We've been recently having some readings related to the fourth \emph{satipaṭṭthāna}, the establishing of mindfulness of specific Dhamma themes. We've started the section on the five hindrances, discussing what they are and how we can work with them in our meditation and in our lives. The five hindrances are sensual desire, ill will, dullness and drowsiness, agitation, restlessness, and doubt. They are embedded habits that rise up when we're in contact with the world. They operate with craving, aversion, or confusion as their root causes and lead to dukkha: they\textquotesingle re unsatisfactory, they're unpleasant, and there is the wish to be free from them.~

These patterns drive our whole life, and we can choose to change them. The Buddha says that through appropriate attention, we can change our perceptual patterns and our experience of the world.

For instance, we don't have to experience the world through the five hindrances. We don't have to seek sensual pleasure in order to have a moment of contentment. We think that engaging in sensual pleasure is enjoyable, and to a certain extent it is, but the actual experience of sensual desire is something I find quite unpleasant. There is an agitating quality to wanting something.

When we compare it to a peaceful and contented mind, the arising of sensual desire is unpleasant. Even if it results in getting something in the sensual realm that's pleasurable, the process is agitating. The goal of sense desire is to get rid of the unpleasant feeling of wanting. And then when we do get some pleasure, it goes away very quickly. All in all, the pursuit of pleasure due to sensual desire is not gratifying in the long term.

All of the hindrances lead to dukkha. The hindrances are unpleasant ways of perceiving and living in the world, default patterns that we've learned and reinforced over a long period of time.

We can start to undo these patterns and reconstruct the ways we perceive the world. We will also become more pleasant to be around. We all know what it's like to be around somebody who's consumed with craving, aversion, or other overt forms of the hindrances. It can be draining to be around them, just as it can be hard to be around ourselves when we're experiencing these hindrances.

Reconstructing these patterns is a gift to the world as well as a gift to ourselves. We need to recognize them first as they arise in our practice and lives. We can't just pretend they're not there or try to paint over or annihilate them. It doesn't work. We have to first experience these patterns. That's the training in the four \emph{satipaṭṭhānas}: to experience the hindrances, know when they're present, know them very deeply and directly. It's not only a conceptual knowing but also allowing ourselves a direct experience of the hindrance. How does it feel in the body? How does it manifest in the mind, the heart? What are its qualities? Is it pleasant or is it painful? Does it cause the arising of suffering? To know experience directly is to know the first Noble Truth.

We have to know and understand the presence of dukkha. After we've seen and validated the experience of a hindrance, we can figure out how to apply an antidote to it. We can observe its arising and passing away. We can reframe our perceptions. We can look at an experience in a different way and turn it around. That's letting go of inappropriate attention and engaging in appropriate attention. For instance, we can use \emph{asubhakammaṭṭhāna} (contemplation of unattractiveness of the body) as an antidote to sensual desire, and we can develop perceptions of goodwill and patience in response to the hindrance of aversion.

We can recreate our responses to the world around us into an experience that's a lot more pleasant to be with. It's still a constructed reality and we'll need to let it unravel completely at some point in the process of deconstructing our way of being in the world. This will lead to a much more sublime, insightful, and penetrating experience leading to the ending of dukkha. The mind that has been repatterned into a more pleasant, open, expansive state of mind allows us to experience the cessation of dukkha much more easily.~

Over the years, I've used contemplation as a way of bringing up skillful perceptions of the opposites of the five hindrances. I pose a question to myself, such as, ``What would it feel like if I were free from sensual desire and I didn't want anything?'' One word that comes up is ``contentment.'' What's it like? We've all had moments of contentment when we don't want anything, when desire isn't arising and we're at peace. The mind isn't surging out towards picking up sensory experience in the realm of sights, sounds, tastes, or tactile sensations. The mind is at rest and~absolutely happy with what is happening right now. Just imagine what that's like, even for a moment. Imagine the feeling of contentment.~

What is it like to be free from ill will or aversion, when we have let go of something that we were gnawing on for a long time? This could be a memory of a conversation or a constructed perception of a person who has given us a hard time. We realize the dukkha in that and let it go or establish an alternative perception. We can create an image that brings a feeling of kindness, gentleness, or open-heartedness. We can tag it with a label like ``\emph{mettā} (loving-kindness)'' or ``gentle.'' What's a good word to describe that quality for yourself?~

One enjoyable memory for me is again about my mother. When I was going through some of the photos of my mother, there was one of her at the nursing home, an image which is very much embedded in my mind. She was experiencing a moderate amount of memory loss and some delirium and was sitting across a table from me. There were a few other people around, including my sister and one of the nurses. My mother was kind of out of it, and we were all talking and engaging with each other. At one point I tried to include her by asking her how she was doing. I wasn't sure whether she was mentally present or not. She didn't look very alert and didn't really engage. I caught her eye and said, ``Mom, are you following all of this? Are you following any of this?'' She just looked a little confused. And then I gently said, ``Do you know who I am?'' She responded with another momentary look of confusion, but then all of a sudden she became alert and was 100 percent present with us. Her big, blue eyes stared out from her wrinkled, ninety-two-year-old face, and with a penetrating gaze, directly into my eyes, she said slowly, ``You are mine.'' And my heart just melted.

In a moment of aversion or going through a rough patch, if I can bring that memory up, it immediately washes that mood away. Instead of aversion, I feel kindness and gentleness. The world opens up again.

We can change our perceptions on the spot by having a word or an image that we can pull up. We can ask ourselves, ``What would it be like to be free from ill will? What would it be like to be free from dullness and drowsiness? What is it like to have a fresh, alert mind?'' You can think of the words ``alert,'' ``fresh.'' What is the feeling that arises? Cool and spacious.~

If you're feeling agitated, ask yourself, ``What's the opposite like? What's it like when you're not agitated? Still, silent, quiet?'' We've all had moments of peace, even if they've been very fleeting. What's the feeling that the word ``stillness'' brings into perceptual experience?~

What is it like to be free from doubt? We feel secure, grounded, stable, unwavering. What does that feel like in the body and in the mind? This self-questioning, then labelling and experiencing the feeling of the antidote, are ways that we reconstruct these perceptual patterns.

When we identify what hindrances are present and experience them directly, we can bring up perceptions of their opposites as a meditation. You can do that as a mantra. After you've established yourself comfortably in the sitting posture, take a few moments to say, ``Okay, what's it like to be contented, gentle, alert? Repeat the words: ``silent,'' ``still,'' ``secure.'' Watch what happens to the heart. It opens, relaxes, and settles. As that happens, the opportunity arises to be present with something simple and soothing, like the breathing. If the mind is not agitated from one of the hindrances, but instead is experiencing positive qualities, this is the mind that composes itself easily.

Freedom from the hindrances can come through reconfiguring our perceptual patterns. It's also important to remind ourselves that these patterns aren't who we are, even the most positive ones. We don't need to take ownership of them and say, ``Finally, I've gotten rid of my self that was full of hindrances, and now I am someone who is full of more positive qualities.'' We're not trying to construct an identity around them either. We're trying to temporarily fabricate skillful states that lead us to a deeper understanding. We create the mind space that can understand the arising and passing of the world around us and how to deconstruct all of our perceptual patterns.

We have to start by recognizing the mind states which obstruct us. We can repattern them in a way that causes us to open and relax and make life enjoyable. We have to enjoy life. Even though we're working towards deconstructing how we view the world and ourselves within it, we have to have patterns that are pleasant to experience. We can then live with ourselves more easily.~

We can be an example for others, too. We become easier to be around if we work through our obstructive patterns and create skillful ones that appear to the world around us. It gives other people hope. We naturally want to live in the world with ease of mind, and this will also help other people.~

The most obstructive perception that we live with is the perception of self. We can't really help other people fully until we have let go of all selfing patterns. We've got to get out of the way if we really want to be of service to the world. If we can learn how to create a deep sense of contentment and happiness in ourselves through eliminating obstructive patterns, particularly the self-perception, then we can be of greater benefit to the world. That's my aspiration anyway. The potential to make these changes is what keeps me here after twenty-five plus years. There is still lots of work to do, but many benefits have come along the way. It gives me hope, faith, and inspiration for the future.~

\textbf{A Sublime Abiding}

Based on a talk given on December 22, 2022

Those of us who have taken monastic vows and those who are committed lay practitioners receive much encouragement to keep in mind the long-term potential of practice: the ultimate goal of realizing liberation. The Buddha talks about it in many ways in the teachings. He instructs us not to become complacent in the practice.

The Dhamma teacher Ruth Denison, who passed away a number of years ago, emphasized the same idea towards the end of her life. A group of people went down from Abhayagiri not long before she passed away to visit her and to wish her well. One of her notable messages was: ``Don't settle for mere happiness. You have the possibility for full liberation.''

The encouragement for realizing the complete goal is a focus of our teachings. We gear our attention towards those teachings that help lead us to the experience of the Unborn, the Uncreated, the Unoriginated, and the Unformed. We keep returning to that goal.

At the same time, we live in the world and have to negotiate our ways around the world, whether it's in the monastery or outside of the monastery. We try to make it a useful place to practice and create skillful, supportive conditions. We have to learn how to negotiate the conditioned realm as part of our path to the Unconditioned.

We are nearing the end of December, and the days are getting darker and shorter. The nights are long. There is a coolness in the air. The signals of nature are moving us towards quiet, turning inwards, contemplation, and stillness.

It's delightful to experience that at this time of year. The holiday messages, at their best, such as those around Christmas, honor the birth and life of a very compassionate figure. The uplifting side of Christmas is remembering qualities of nurturing, care, concern, compassion, and skillfulness in our relationships. There is more attention on family, friends, and the people with whom we live, and on the goodness that can be cultivated from these relationships.~

It's also easy to get lost in some of the less desirable aspects of the season. There is the side of consumer mentality and acquisitiveness, the material world. Businesses make a lot of money at this time of year and can get caught up in seeking gain by selling their products; others can get very absorbed in the consumer ethic.~

Going into a shopping mall at this time of year, with gaudy decorations and blaring music, can be jolting. For those of us who are living in or around the monastery and lead a quiet life, it's an affront to the senses. It's a mixed time of year, depending on what we focus on.~

Learning how to live skillfully in the conditioned realm is part of the path to the Unconditioned. This means learning how to live in the world, but not be of the world. We learn how not to get lost in the worldliness of the world and how to navigate in a way that is still conducive to the path.

Even in the monastery, we don't isolate ourselves and shun the world. We do try to step away from the worldliness of unskillful actions and nurture skillful ways of being in the world, but we're not apart from it. One of the most skillful ways of developing ourselves in the realm of the conditioned is the development and abiding in the four \emph{brahmavihāras}, the divine abidings.

\emph{Vihāra} means dwelling place, so brahmavihāra means dwelling or abiding in the mind state of a brahma god. These are beautiful states of mind that beings in high spiritual realms dwell in for long periods of time. We can also cultivate them and make them an abiding in our lives in this world. Even though these states are of the conditioned realm, they are a very important part of the path to total freedom.~

The four brahmavihāras are \emph{mettā} (loving-kindness), \emph{karuṇā} (compassion), \emph{muditā} (sympathetic joy), and \emph{upekkhā} (equanimity). Many of the current teachings regarding the brahmavihāras focus on specific meditation techniques to develop them. The emphasis is on an active practice of \emph{mettā-bhāvanā} using various mental recitations, verbal phrases, such as, ``May I be well and happy.'' There certainly is value in this approach for many people. For others, including myself, an alternative approach to developing the divine abidings is first investigating what mind state might be obstructing the natural arising of the brahmavihāras. We look closely at the \emph{kilesas}, the obstructive mind states that stand in the way of opening to sublime mind states.~

If we don't clearly see the mental obstructions that block the brahmavihāras, then we're possibly only doing it halfway. If we don't notice the underlying tendencies that lurk as potentials, and if we don't see them clearly, then it\textquotesingle s so easy for them to catch us off guard, even if we don't think we have such states of mind very often. If we don't attend to them with clarity and honesty, then they may rise up unexpectedly, and take us by surprise. As with the First Noble Truth, we have to fully understand the experience of dukkha before we can start to let it go.

Obstructive mind states are stimulated by outside circumstances and we have to meet them with honesty and clear awareness, without getting lost in them and without moving into and indulging in them.

For the first brahmavihāra of mettā, the obstructive states that can rear up are ill will, anger, aversion, and all of their unfortunate manifestations. Somebody who is unskillful or unpleasant to be around may trigger something in us that rises up as ill will, aversion, or anger. We shouldn't develop aversion to the aversion and try to turn away from it, pretending it's not there. We shouldn't bury it or feel like we shouldn't feel that way. We should turn gently towards the aversion to fully understand it, to fully face it and see that state of mind as it is experienced right in the present moment. This is done most skillfully by attending to the bodily energies.

If we get involved in the storyline too much, ``He did this. She said that,'' then we're just increasing the proliferative tendency. We need to drop the story and feel the feeling. We need to turn to the affective tone of the experience to really understand it without getting lost in proliferation. We should develop the bravery to face the unpleasant feeling. There may be clenching in the stomach, tightness in the cheeks, contraction of the facial muscles, or the rising of heat. Sometimes, if the feeling is quite intense, there's a flush that's really unpleasant, and our normal reaction is to try and get away from it. Nobody wants to have unpleasant feelings. In order not to deal with the pain, we can get lost in the storyline, trying to justify, explain, or present our counter-arguments, with all kinds of conversations in our head.

With patience, though, we can get past that discomfort. We can bear up and experience the pain by facing it and allowing it to manifest as an energetic quality through the body. We can gently bring our breathing, kind attention, and spaciousness to the experience and not get wrapped up in the storyline. We can attend to the contracted state on the bodily level and then more deeply in the heart. We allow ourselves the space to be fully with the experience, without absorbing into and encouraging it.~

Fully attending to intense bodily energies takes a lot of patience and spacious kindness and further encourages the qualities of kindness, spaciousness, gentleness, and warmth to arise. It doesn't mean that we're endorsing unskillful behavior, but we're learning to experience a broader sense of kindness, including well-wishing to ourselves, by owning up to the truth that people make mistakes. We make mistakes and other people make mistakes. We can disentangle ourselves from the enmeshed obstructive mind state and wish ourselves and others freedom from those reactions.

We can realize that people who are triggering our aversion, anger, and ill will, are probably feeling the same. We can wish them well. If they can experience relaxation from tense mind states, then they will act in ways that are much more pleasant for us, too. It's a win-win situation.~

The afflictive state of wishing harm is the opposite of karuṇā or compassion, the second brahmavihāra. We might see this arise in subtle ways when somebody who has been unskillful receives the karmic consequence of their action and experiences some misfortune as a result. When the tables turn and they get their comeuppance, we might feel a momentary, brief delight that they are ``getting theirs'' now. This happens a lot in the world of politics, where there is so much anger, tension, and contention. A wish for revenge emerges when the tides turn, which they always do.

There may be some insensitivity or cruelty when somebody is experiencing suffering and difficulty. There may be an inclination, for whatever reason, to turn a cold shoulder to that person. That's a really unpleasant state of mind. If that happens, even on a small scale, what does that feel like? We need to turn our attention towards the unpleasant feeling of constriction and tightness when we experience moments of cool indifference or insensitivity towards somebody else's suffering.

By experiencing the results of our own insensitivity, we realize that other people are no different than we are. Compassion means learning to be with the passions and difficulties as they manifest and recognizing that somebody else's suffering is the same as our suffering.

Compassion is the quality of quivering or trembling with, to be in resonance with someone else's suffering. When we know our own suffering, then the natural response to another's suffering will be to feel it at some level. As much as we want to be free from suffering, we want others to be free from suffering. It's all the same.~

The third brahmavihāra is sympathetic joy. An afflictive quality that we can experience related to the third brahmavihāra is envy when someone else is experiencing success. Envy is an unpleasant state of mind that arises when we want to get someone's attention, care, or concern, but they are giving it to somebody else. As they are giving someone support or offering others praise and recognition, we're thinking, ``Well, what about me? What about \emph{me}?'' We ignore the fact that somebody else is receiving justifiable recognition or support and that they're feeling good. We get lost in our ``what about me?''~

When that envy arises, we can respond with the same approach as with aversion. We fully experience the bodily energy and say to ourselves, ``Okay, well, this is some sort of conditioned wish for recognition.'' As long as we are subject to praise and blame and success and failure, when somebody else gets something we would like, envy can arise. We need to know this honestly and not get lost in self-flagellation for experiencing afflictive states because they are part of our conditioning, part of the patterns that we have developed for whatever reason. We need to recognize them as energetic states and then allow them to pass. If we attend to the energetic state with kind attention, it will dissipate.~

Once the negative state passes, we can turn towards that other being who is receiving something pleasant and recognize that when we receive that kind of attention, we also feel pleasure. To resonate with another person's success is a rare commodity in the world. To really open up and feel genuine delight when someone else has received something is quite wonderful. That is muditā, the brahmavihāra of sympathetic joy.

The last brahmavihāra is equanimity. When we're not equanimous we are pulled around, as with the eight worldly winds. We get pulled towards pleasure and away from pain. We are buffeted around by seeking gain and avoiding loss. We are affected by praise and blame and by seeking good reputation and avoiding disrepute. We are consumed with bouncing back and forth in reaction to the ways of the world, the world of opposites. Realizing the way the world is, we can choose either to get lost in it, repeatedly being overwhelmed by opposing states of mind, or we can choose to see them with clarity.~We can develop the ability to be with these opposites with a clear, knowing state of mind. Equanimity has a very strong quality of discernment, which we can employ in our encounters with the world. We can learn to step back from over-reacting to conditions because they are always going to fluctuate. They will always change.~

Our basic practice mode is to own up to the underlying states that obstruct the brahmavihāras. We then encourage the states of kindness, sensitivity, compassion, warmth, well-wishing, and quiet joy. We experience our own gains and successes and those of others, without attachment. We can experience evenness and an understanding of the ways of the world. These qualities are dwelling places in our hearts where we can settle, not just as a practice, but as an actual abiding. We steadily make that our frame of reference. By seeing the disadvantage of how we normally operate, we step away from the normal ways of the world. We then can live in the heavenly realms in our minds, allowing that space to open up and develop.~

Equanimity is so important in and of itself, and it's an important part of the other three brahmavihāras, too. When we have equanimity, we are less likely to get lost in what the commentaries call the near enemies of the brahmavihāras. Equanimity provides an alert clarity and awareness of our states of mind. In regards to the development of mettā, for example, equanimity prevents us from sliding into attachment, which is the near enemy of mettā. With clarity, we move beyond the personal fondness for one specific person and cultivate a global quality of kindness and well-wishing, independent of individuality. Mettā is a feeling for all beings and a very sublime state to abide in. Attachment narrows our attention down to one individual or a limited number of people that we happen to like or are attracted to.~

When we are developing loving-kindness reflections, we shouldn't let our minds slide into affection for an individual or individuals. Even if we are not moving into sexual desire, Mara is oftentimes waiting on the edges for an opportunity to cloud the mind. The wonderful quality of loving-kindness for all beings can turn into something more involved with attachment, clinging, and confusion. Equanimity and wisdom help to prevent us from sliding into the near enemy of mettā.~

The same is true with karuṇā, compassion. When we are feeling compassion for someone and their suffering, if we are not careful, not mindful, not equanimous, we can slide into a feeling of overwhelm. We become identified with the suffering of someone else, and we don't really know how to handle it. We absorb into it and don't have the proper boundaries that come from clarity. Whether it's our own or someone else's, we get overwhelmed with the suffering of the world. Equanimity keeps compassion stable, clear, and balanced, so we know when and how we can help without getting lost in emotion or reactivity.

The brahmavihāra of muditā can tend towards overexcitement, even to the point of agitation about something good. When a family member, partner, or monastic friend is experiencing something wonderful and we get too excited, the sense of balance disappears. We've lost the perspective of ``well, yeah, this is good, and to be delighted in, but it's also part of the conditioned realm.''~

To keep even a good thing going takes constant management, and things will probably shift. One may not always experience this kind of success and the tides might turn. The blame might happen again. The loss might happen again. This isn't meant to negate the good quality of the experience or to look at it as sour grapes: ``Oh, it's going to go away anyway.'' We can delight in goodness \emph{and} also reflect that it's a conditioned experience that will change and fade over time.~

Equanimity enhances and stabilizes the other three brahmavihāras. However, equanimity itself can sink into its near enemy of indifference. Sometimes we can slide into thinking, ``Oh, well, everything's impermanent, so what does it matter?'' If everything's impermanent and it's going to change, does that lead us to indifference or even nihilism? If we let equanimity slide in that direction, we lose the benefits of all of the other brahmavihāras. We won't be able to sustain a sense of well-wishing (the wish for the end of suffering) or delight in joyful experiences.

Indifference is not at all what equanimity is. Equanimity involves clear seeing. The states of mind that we are trying to develop are the results of intentional skillful reflection and good actions. They are the results of skillful kamma. Equanimity includes the wisdom that knows how these states come into being and how we can maintain them as skillful qualities. This is Right Effort: raising up, establishing, and maintaining the wholesome qualities that serve as parts of the path.

As we develop the four brahmavihāras, we also realize that we can develop them to the point where they are ``abundant, exalted, and immeasurable.'' These are very expansive, spacious states of mind. The process of developing them can result in exalted, abundant, and measureless states of mind. By letting go of the constricted states of mind, the mind opens, relaxes, expands, and broadens. The experience of each of the individual four brahmavihāras grows into a very wide, open state of awareness.

But it all comes back to recognizing and owning up to the \emph{kilesas}, the constricted states of mind, and letting them go. This is done through clear understanding, not by avoiding them, and by allowing ourselves to experience them. As we let go of these underlying states, we can open up to the comfort and happiness that arises.

The brahmavihāras are the best we can do in the conditioned realm. Even on our quest towards realizing the Unconditioned, we have to negotiate our way through the world and react in the most skillful way to the conditions around us. The more we do that, the more this light ease of being can manifest. The old habits slowly wear away. While we are still part of the world, we develop a more beautiful way to be in the world. The world becomes a place where we can comfortably abide as we negotiate closer to realizing the Unconditioned.

\textbf{The Delusion of Forever}~

Based on a talk given on February 10, 2013~

~

Luang Por Pasanno's most recent teachings during this Winter Retreat have been exploring the five \emph{khandha}s, the five aggregates. The term aggregates refers to five categories of personal experience: form (the body), feeling, perception, volitional formations, and consciousness. These are the aggregates of existence.~~

~

The Buddha doesn't mince words in his teachings and with his straightforwardness and directness; he uses some pretty bold expressions, especially in reference to the khandhas.~One of the phrases that often occurs in the teachings refers to the five aggregates as a disease, a dart, a tumor, a boil, a calamity. These are strong words. They point to the potential for the aggregates to cause suffering when we invest them with a sense of ownership and use them as a means for seeking satisfaction in the world.~

~

One of the teachings from the Saṃyutta Nikāya is titled ``Yamaka'' (SN 22.85), the name of one of the monks.~To summarize, Yamaka is possessed of an annihilationist view and the monks try to dispossess him of this view. They are not successful and so they consult with Sāriputta, a chief disciple of the Buddha, who does an incredible job of pointing out his wrong view and directing him to understand that the Buddha's teaching is neither annihilationist nor eternalist.

Sāriputta then uses a simile that is very powerful.~The story relates there is a wealthy householder, and there is a man who has something against him and wants to murder the wealthy man. He thinks to himself, ``Suppose I become his servant, get to know him, and work my way into his life.'' He approaches the wealthy man and offers his services as a servant. The householder accepts this offer, allows him to come into his life, and becomes familiar with him. Over time, with diligence, the servant attends to him, providing him with all he needs and gets to know the householder well. He rises in the morning before the householder, ready to greet him, and stays up until after he has gone to sleep. He is attentive and endearing in his speech and actions, slowly gaining the wealthy man's trust.~

After some time, with a comfortable routine well-established, the householder lets his guard down completely. The servant then seizes the opportunity and stabs him in the back, killing him.

Sāriputta then asks Yamaka, ``When this servant was initially offering his services, was he a murderer then, even though the householder didn't know it?'' Yamaka replies, ``Yes.''~

``And when he was getting to know him, serving him, and setting the stage, was he a murderer then?'' ``Yes.''~

``And when he finally stabbed him in the back when the householder wasn't looking, was he a murderer then?'' Yamaka answers, ``Yes, he was.'' So even though he didn't know it, there was a murderer in the house the entire time.

~

The murderer in this story is a metaphor for the five khandhas. We place a lot of trust in them, depend on them, and think they are here for us. They serve us well and provide us with pleasurable, warm experiences, and we become lulled by the pleasure they provide. All the while, we are not paying attention to the fact that at some point, they can and will turn against us. We willingly buy into this paradigm and believe in it.~The khandas are there for us when we go to bed. We sleep, dream, wake up and the khandhas are still there, ready to offer us what we want to gratify our desires.~

We know that the teachings tell us not to put our faith in the aggregates, yet we still do. Time after time, we get pulled into the belief that we have a right to expect the khandhas to operate the way we want them to operate; that they should do what we want them to do. In the United States, we grow up with a constitutional right to life, liberty, and the pursuit of happiness, and we incorporate that into our belief system. And then something goes wrong.

Many of us are having physical health problems. Various muscles, joints, and internal organs are injured or deteriorating. The body doesn't obey us the way we think it should. The Buddha gave the monastics permission to use the four requisites, one of which is medicine for when we are sick. We have access to medical care, and we do what we can to keep the body relatively healthy.~

The problem comes when we expect it to stay that way. We get a treatment, and maybe it doesn't work. We have an operation, and it fails. Perhaps the problem isn't what we thought it was, and it's going to take longer to fix than we think, or it's not going to get better at all.~

~

One time, before I went in for surgery to repair a hernia, I thought, ``Oh, I'll be up and moving comfortably in a couple of days.''~For the surgery, I was given an anesthetic and morphine. I went to sleep and woke up a little while later. When I was more alert, the doctor came over to me and said, ``Well, it was bigger than we were anticipating, so you're going to be in recovery for a longer time.'' I had blissfully been put to sleep, woke up, and found out that I had been duped by this body.~

It happens to all of us in different ways.~We recover from an illness and we think, ``Now I'll be able to really get down to practice. My energy will return and everything will be fine.'' We believe the problem is cured and there is not going to be any more difficulty; we're home free. Maybe that assumption lasts for a little while, but then something else happens.~

It's not that we don't want to keep ourselves healthy and keep our bodies going as best we can. But if there is any expectation that it will be smooth sailing from then on, we are setting ourselves up for the next round of suffering. It is like putting a Band-Aid on a large, gaping wound, thinking it will heal it forever, or like plugging the hole in the Titanic with a piece of chewing gum. It's really a desperate fantasy.~~

~

Many years ago, a longtime friend of the monastery was hiking and camping in the mountains. One evening he lay down in a field to sleep and woke up a quadriplegic.~He had been run over by a vehicle that was traveling across the field. He wasn't expecting that when he went to sleep that night. This friend has reflected that the rest of us with ``able bodies'' refer to people in his situation as ``disabled,''~but sometimes people we call ``disabled'' refer to us as TABs, ``temporarily able-bodied.'' That is a reminder that no condition is permanent.

~

Even the body itself inclines towards renewal, health, and revitalization, sometimes in the face of great odds. A strong image that remains in my mind is from when our neighbor was diagnosed with colon cancer and was starting to realize that it would be the cause of her demise.~The cancer had spread to many places in her body, and she had surgery to remove part of the diseased area.~At that time, she was clearly in the last few weeks of her life. I remember her lying there with her digestive system shutting down. She wasn't taking very much nourishment and was starting to waste away. Other organ systems were~also shutting down.~Even though she was getting IV fluids, the kidneys and circulatory system weren't functioning properly. The fluid~was collecting in her legs and lungs, and at times she was having trouble breathing. We could see the body was packing it in.

One time when the visiting nurse was changing the dressing on one of the surgical incisions from the operation, it was clear the wound was healing, it was getting better. Everybody was commenting, ``Oh, that looks really good. That's not going to cause you any more problems. We won't have to change the bandage anymore. It's healing up.'' It hit me as being so bizarre that she didn't have a lot longer to live---her body's major organ systems were all failing---but still, the body was working to heal itself, trying to mend this little wound.~It struck me as so paradoxical.~

The body itself seems hardwired to support us in this delusion of permanence.~Then, when we turn our backs, it proves otherwise. It's the same with the other khandhas: feelings, perceptions, mental formations, and consciousness.~We have an expectation, a hope, a fantasy that as soon as we get rid of an unpleasant feeling and find a way to replace it with a pleasant feeling, everything will be okay. So we search for something, either in the material realm or the mental realm, to provide a pleasant feeling. If we obtain it, we relax and think, ``Ah, this is the way it's supposed to be.'' This can even be the pleasantness of a ``good'' meditation. We rate our meditations as a ``bad'' meditation or a ``good'' meditation based on the level of bliss that we experience. We come to expect that a pleasant meditation is what we should have. Then, something changes and pleasant turns to unpleasant.

Regarding the aggregate of perception, we can cling to a perception so strongly that when something comes along to disprove it, it is very hard to change and perceive it differently. The Buddha refers to perceptions as mirages because it is so easy to view people and events through the lens of distorted impressions that are not accurate representations of those people and events.~

Think of people you know. In terms of personal experience, people are just a perception in the mind---how easily that changes. Someone you have a lot of good times with can suddenly change and do something that takes you aback.~Your perception changes dramatically and you experience him or her as a completely different person. All of a sudden, you find out that your faithful partner of twenty years is shacking up with the secretary or inviting the mailman in for lunch.~What happened to that wonderful person in your life, who was always so warm and supportive and always there when you needed him or her?

More commonly, someone that you get along with has a bad day and says something sharp or angry, and all of a sudden he or she is an evil monster, at least for a period of time.~Just when you think you have got it all worked out with the people around us and everybody is harmonizing, something happens and it changes. You're caught by surprise.~~

~

The fourth aggregate of mental formations includes our moods and emotions: how much control do we really have over them? We work with them and try to be more skillful by cultivating wholesome states of mind.~That is part of the path. But when we think that they should be one way and that they shouldn't be another way, then we entangle ourselves even more.~We create another mood or thought that is unpleasant, negative, or averse. We can slip into further despair in reaction to that negative mind state. Once again, we haven't lived up to our expectations, and we blame ourselves, compounding the mind state with another layer of suffering.~

~

Consciousness, the fifth aggregate, is constantly changing. The nature of consciousness is to weave the illusion of continuity. The continuity of experience through consciousness gives us a sense of solidity and identity.~Consciousness constantly pulls data from the six senses. For example, visual consciousness is ignited by an interesting object, but boredom quickly sets in. It is interesting for a flash and then the desire to experience the object fades just as quickly. An interesting sound arises, pleasant or unpleasant, and captivates our attention.~We pull that into the mind and suddenly we are experiencing ear consciousness. The same is true with the other sense bases: smells, tastes, tactile sensations, and ideas of the mind. Attention is directed toward an object in one sense base, abandons it and immediately takes in another object, perhaps from a different sense base. We weave this together with the perception of ``this is who I am.''~All of these moments coming together in rapid succession create a sense of continuity, like a movie being projected onto a screen, and we believe in it.~

This is how the five khandhas operate when they're affected by craving, clinging, and delusion. They keep us feeling secure, entertained, and feeling like this is who we are. We constantly renew this perception over and over again, but unexpectedly, the khandhas betray us.~

However, there is also an interesting paradox in the teachings. In order to free ourselves from this confusion, the teachings point us back to the five khandhas. They are not only our ``murderer,'' but also the way out of this whole mass of suffering.~~

~

There is a sutta where a young deva of stunning beauty appears in front of the Buddha, illuminating Jeta's Grove (SN 2.26).~His name is Rohitassa; he pays respects to the Buddha and asks him a question. He says, ``Bhante, is it possible by traveling that one can reach the end of the world where there is no birth, no aging, no death, no passing away, and no rebirth?'' The Buddha replies, ``No, Rohitassa, there is no way by traveling that one can reach that end of the world where there is no birth, aging, death, or rebirth.''~Rohitassa says, ``Oh, it is wonderful that you say that because in a past life I was a seer, a sage, referred to as a `skywalker.' With one stride, I could step across an ocean, and my speed was rapid like that of a light, swift arrow soaring through the sky. I had a lifespan of a hundred years, and I spent my life coursing through the cosmos, with my quick and long stride, trying to find the end of the world where there is no birth, aging, death, and rebirth. But I died before I got to the end of the world.''~The Buddha responds, ``Yes, there is no way that by traveling you can reach the end of the world, the end of suffering. I tell you that it's just in this fathom-long body where there is the world, the origin of the world, the cessation of the world, and the way leading to the cessation of the world.''~~

The path to the end of suffering is found by exploring the experience within our bodies. The examination of these five khandhas, of what they really are and are not, is the way leading to the end of suffering.

We have to be willing to throw our whole life into our spiritual practice, every bit of energy we have. We have to be willing to make that commitment because the habits are strong. The tendency to want to find satisfaction through the experience of these five khandhas is very strong.~At the same time, the full-on commitment has to be made with compassion, kindness, and gentleness. Commitment is not something that can be forced on us. It has to be taken up with much care, attention, and clarity.~

~

It is through investigation of the khandhas that we learn to know what has been fooling us all along. We don't just disregard the need for this clear and complete examination. We don't run away from it. We turn around and look carefully at the khandhas. We turn around and face the murderer, which is not easy to do. We look at all that we have been investing in for so long, with fresh eyes. We need to bring a lot of mindfulness and alertness to the process.~

~

At the same time, we have to be willing to recondition our perceptions.~We can cultivate skillful mind states that result in more useful perceptions. The new perceptions are also subject to change and deterioration, but we use them in a skillful way. We use mindfulness and alertness to develop all the path factors, including renunciation, kindness, and gentleness. We use the aggregates that have been fooling us, turn them into instruments for a deeper understanding, and see them in their true nature. We see through the illusion of permanence and solidity and realize that the aggregates cannot be a refuge. We can, however, use them in a skillful way in order to understand how they really operate.~

~

When we accurately see how the khandhas operate, over and over again, disenchantment with searching for any permanent sense of satisfaction from the aggregates arises, and identification with them weakens. We have avoided seeing the aggregates clearly for so long and now we have the opportunity to look at them squarely in the face. We use them for our long-term happiness and ultimate liberation.~

\textbf{Turning Toward a True Refuge}

Based on a talk given on November 27, 2020

This morning, we received a teaching on the contemplation of the five khandhas, and in particular, the \emph{upādānakkhandhas,} the aggregates affected by \emph{upādāna} (clinging). The reflection included the kinds of clinging, different descriptions of the khandhas, similes that the Buddha used, and the importance of contemplating them through the lenses of \emph{anicca} (impermanence), \emph{dukkha} (suffering), and \emph{anatta} (not-self).

Contemplation of the aggregates should be one of our regular meditation practices as it's at the crux of the path to freedom. It deepens our understanding of what we usually regard as me or mine. Developing insight into the process of I-making and mine-making will lead to seeing the bald-faced reality of what it is that we have been taking refuge in. This understanding leads to a sense of disenchantment and supports relinquishment. This is the path to freedom. Freedom is not found in becoming anything special or having all sorts of high consciousness experiences. It comes from seeing through our presumptions of stability, close examination of where we seek pleasure, and most importantly, contemplation of the perception of self. We look closely at what we are taking refuge in, questioning the validity of that refuge as a long-term source of happiness. We need to see the truth of our reality to the point where we want to turn away from that false refuge and see something deeper.

Most of us have been doing this routine of pleasure-seeking long enough to have a sense that it's not taking us exactly where we want to go. This realization is the impetus, the inspiration for seeking something beyond this experience of me trying to be happy, me trying to find satisfaction in the world. We are looking to actualize the transcendent goal of freedom instead of constantly striving for temporary happiness.

In the suttas, there is a discussion between Venerable Sāriputta and Venerable Mahākoṭṭhita that emphasizes reflection on the five khandhas (SN 22.122). The discussion of what to look for in examining the five khandhas is found in many different suttas and includes taking stock of the five khandhas, form, feeling, perceptions, mental formations, and consciousness, and experiencing them through the lenses of anicca, dukkha, and anatta.

Several other words are used to describe how to regard the five khandhas, and they have a powerful impact when we hear them recited. We view each of the five khandhas in turn, starting with form, as impermanent, suffering, a dart, a tumor, an affliction, an illness, a misery, a disintegration, alien, empty, and not-self. These are very strong words. Similarly, when we experience and examine the khandhas of feeling, perception, mental formations, or consciousness, we see them all in the light of impermanence, suffering, a tumor, a dart, an illness, an affliction, a misery, disintegrating, empty, and not-self. These words grab our attention and give rise to alert reflection.

As a result of reflecting on these words, we can ask, ``What am I doing taking refuge in the five khandhas?'' Disenchantment arises from the realization of where we have put all of our hopes and attention.

We can start to reflect deeply on the five khandhas, to examine them and to see if we can uproot the perception of self in relation to them. The perception of self comes in different forms: as ``I am,'' ``I am this body,'' or ``I am these feelings, perceptions, thoughts, or consciousness.'' The notion of self can also manifest as a perception of ownership, the experience that they are mine. ``This is my body. There is somebody here who owns this body, who possesses these feelings, perceptions, thoughts, or consciousness.'' We may not view them as ``me'' or ``I am some part of these or all of these,'' but there still can be a subtle view of somebody who is here, who possesses them, owns them and should have some control over them. That's what a self would do. That's what a ``me'' would do if it were solid, real, and in control.

We need to pick apart self-view, experience it, and question it closely, which is what our settling and calming meditation is designed to support. We aim to get clear, steady, contented, and peaceful enough to be able to hold a contemplation like this, to be with it for a while and delve into it. It doesn't mean we can't think about it in a skillful way before our minds collect and pull together. In fact, engaging in active contemplations can help settle the mind even more. We talk about \emph{samatha} (calm) and \emph{vipassanā} (insight) being two ends of the same stick. They help develop and support each other.

Contemplation of the body encourages exploring the body in relation to anicca, dukkha, and anatta and leads to questioning the assumptions about the nature of the body and our firmly entrenched ways of viewing ourselves. There are many different ways that we can contemplate the body, and it helps to start with a firm determination to make this our Dhamma lens, our focus for exploration.

Earlier this year (2020), when I stepped down from being co-abbot with Ajahn Nyāṇiko, he put together a lovely book of photographs from throughout my life and personal recollections from many people, mostly from my monastic life. It's a great collection. One of the photographs he got from my sister was of me when I was less than two years old, a photo I had never seen before. My mother was holding me in her right arm, balanced off to one side, and my brother, sister, grandmother, and a young neighbor friend were in the photo. I had a slight scowl on my face and looked like I wanted to get down. When I look at that photo, the question comes up, ``Was that me?'' It doesn't seem like me now. It was a touching image, but it brings up the question, ``Why would I call that me, when it feels like this me is so different from that me?''

There was another photo of the Sanghapāla Board from the time before I was ordained. It was before Abhayagiri Monastery was formed and was during one of Ajahn Amaro's trips to the US, when he stayed for a few months. Towards the end of this visit, there was a Sanghapāla Board meeting, and although I wasn't on the Sanghapāla Board, I came down from Seattle to join the meeting. I was in my late 30s, as were Debbie and Ajahn Amaro, who were both in the photograph. The three of us were sitting there looking quite young and thin. As I was looking at that photograph, I thought, ``Well, Ajahn Amaro looks pretty much the same. He has aged some and put on a little bit more weight, but Debbie and I look substantially different.'' In the photo, I had brown hair, a beard, and was wearing trousers and a sweater. When I look in the mirror now, I see almost no hair, certainly no beard, and what hair is present is white instead of brown. I'm clearly not in my late 30s anymore.

The photos were of when I was two years old and in my late 30s, but I am now in my mid 60s. Is this the same person? Biologically, we know that human body cells change out more or less every seven years, some much faster. Some turn over on an almost daily basis, and others are quite long lasting. Anatomically, we're not the same being that we were years ago. Where did that ``me'' go? Most of the cells that were me, that made up this body seven years ago, have died and have been excreted through urine or feces. They have been replaced with similar ones, and the body in general has morphed and aged. There is merely an illusion of continuity that carries through.

The same is true for all feelings, perceptions, and memories. Where did Tom DeMaria go? Where did Sāmaṇera Karuṇadhammo go? These are all just creations in our minds, even the Karuṇadhammo Bhikkhu of yesterday, who was sitting in this spot answering questions. Where did he go? It's just a memory. Everything prior to this moment is just a memory. The same is true for everyone. We can reflect on our assumptions and views of who we are and realize that this perception of self is a manufactured illusion to keep a sense of continuity in our lives. But it's a construct. It's an assumed reality. Realizing this truth can be quite a relief.

It's easy and so tempting to hang on to memories and make them into present realities. Many memories are painful, and hanging on to them leads to more and more dukkha. Instead, it is a relief to be able to say, ``Oh, that event is finished and is only a memory. I don't need to impute any particular reality on to it other than as a memory.'' And then here I am now. Here, just as it is. The future is a guess, a projection. There really isn't anything happening other than what's happening right here and now.

Reflecting on the illusion of continuity is one way we can start to get a handle on the perception or assumption of self. Another way to view it is by looking through the lens of control. We can investigate what we have control over and what we don't have control over, and how that affects self-view.

There is a modicum of control relative to our bodies and minds. Intentions arise to perform an action through either body, speech, or mind. An intention arises based on some cause or condition, such as sensory contact or a memory of an event. Following the moment of contact, an underlying tendency of desire or aversion gives rise to an intention to do something, to act or speak.

But is there really an independent agent, a self, who is in control? The experience of intention arising is often filled with the notion of me or mine, which imbues the action with a perception of control. We do have the ability to control our experience to a limited extent. For example, I can form the intention to raise my right hand, and then my hand goes up. That means I'm in control of my hand, right? I do what I want to do with it when I want to do it, but it only goes so far. If we reflect on the body, we notice it follows a natural course of experiences which we don't have any control over. The aging process keeps on going. We can say to ourselves, ``Okay, body, don't age!'' but it doesn't pay any attention to us. The intention falls flat. There is no control in this case. Old age, sickness, and death are inevitable. We can influence our health, modify it with the food we eat, the medicines we take, or the exercises we engage in. But in reality, we're all in different places on the same conveyor belt, and it's moving in one direction. We need to closely investigate the sense of control. Who's in charge of this body? Nature is in charge. In quiet moments, we can have this insight when we deeply investigate the notion of control relative to the body.

Right now, sitting quietly, compose the mind and tune into the presence of the body. Allow the body to relax, and keep your focus on the whole body, not letting your attention wander around. Start to notice the various sensations, feelings, and sensitivity of the body. A phrase that one of our teachers uses for this bodily presence is ``intelligent sensitivity.'' This body is one of intelligent sensitivity. It has its own form of intelligence and knows what it needs to do. It will mostly keep on doing what it needs to do in the best way it can, oftentimes in opposition of what we intend. You can notice that when you are sitting still. Notice what are called autonomic functions of the body. The autonomic nervous system is that which controls the breathing, heartbeat, digestion, and other bodily processes. These keep on functioning, regardless of our input and wishes. You think you should have control, but tune into the breathing process for a moment and ask yourself how much control you have. You can command yourself to hold your breath for a short while, but how long does that last? It can be interesting to hold your breath for as long as you can with the intention of having control. After ten or twenty seconds, or a little bit longer, the body overrides any intention that you have and will gasp and take a breath. It's just what bodies do.

Alternatively, tune into your heartbeat. You can feel the heart rate speed up a little bit through intense mental activity or when you are excited or angry. You might be able to slow it down consciously during meditation. However, the heart keeps beating sixty to eighty times a minute, every hour, twenty-four hours a day, seven days a week, 365 days a year, for however many years you live. The heart has an intelligence that is not a cognitive intelligence, but a bodily intelligence. It knows what to do whether you're paying attention to it or not.

You can also pay attention to other sensations, such as slight muscle movements. You're probably shifting position frequently in some way or another, maybe in a very small way. You're not consciously intending every movement, the body is just doing it. Be with the processes of subtle, twitching muscles of the face, or of swallowing saliva. Scan through the body and sense the presence of the body and its movements, its energies, its elements. Gently ask yourself, ``Who's in charge?'' In actuality, no one is. What a relief!

You can reflect on the nature of control in regard to the other aggregates. For example, you can contemplate the pleasant and unpleasant feelings that arise and fade away. The Buddha compared feelings to bubbles forming and bursting in a water puddle when it rains. Feelings arise and pass away, one after another, after another, after another. Sense impressions, memories, and perceptions all do the same thing. Likewise with thoughts and emotions: Where did that mood go that you were wrapped up in yesterday or the day before? Those strings of stories and thoughts, where are they now? Do you really have control or are the defilements producing them? Greed, anger, irritation, and the different forms of confusion or anxiety are generating thoughts and feelings. There's not anybody in there doing it.

There is also the good stuff. By putting attention on developing generosity, kindness, patience, and other skillful mental qualities, pleasant wholesome thoughts, feelings, and emotions will result. But, again, it's not ``me'' who's doing it.

Lastly, the consciousness aggregate weaves together its story of our world by jumping from one sense base to another and creating a story, but there is no independent, self-existent being doing it. This aspect of experience is only a series of perceptions and feelings, which lead to intentions, thoughts, or bodily actions, both skillful and unskillful, which circle around to produce more perceptions and feelings.

We can use the quiet, contemplative mind to question who we are, what we assume about our reality, and what we assume about ourselves. It is all just an assumed and constructed reality that seems very real when we're engrossed in it, involved with it, and acting on it. But when we take time to look and see these mental activities just for what they are, the fascination with what we can get from them lessens. We can gratify a desire or a craving, such as a craving for being or existence, but these are just fleeting desires, finding expression through the five khandhas. We can use the power of a calmed and settled mind for strategic questioning of who we think we are.

Similar to the discussion regarding how to view the five aggregates mentioned earlier between Sāriputta and Mahākoṭṭhita, there is a sutta in the Aṅguttara Nikāya in which the Buddha instructs a disciple in a similar way (AN 9.36). While engaged in the first \emph{jhāna}, the disciple is to regard any phenomena that arise in relation to the five khandhas as impermanent, suffering, a dart, a tumor, an affliction, a misery, disintegration, alien, empty, not-self. The Buddha then goes further and instructs the disciple to turn away from those phenomena and instead direct the mind to the deathless element. The instruction to ``direct one's mind to the deathless element'' occurs in several places in the teachings and does not always include an explanation (for example, AN 3.131). How does one do that? If one hasn't experienced the deathless element directly, how does one know how to direct one's attention to it? But this particular sutta in the Aṅguttara Nikāya continues with, ``Direct one's mind to the deathless element thus: `This is peaceful, this is sublime, that is, the stilling of all formations, the relinquishment of all acquisitions, the destruction of craving, dispassion, cessation, Nibbāna.'\,'' This is a description of how to directly contemplate the deathless element. But if the mind is not yet ready for the direct experience, this instruction can be used as a perceptual training to imagine what the peace of Nibbāna would be like. Simply reciting these phrases can be a settling and insightful contemplation.

With the dispassion that results from seeing the khandhas as impermanent, unsatisfactory, and not-self, we are no longer enthralled with the temporary gratification that the khandhas offer. We can incline towards reflecting on a much deeper peace, the peace of Nibbāna. This is peaceful, this is sublime. The stilling of all formations. That's when the kammic process comes to a standstill. We are not creating any new kamma. Formations (saṅkhāras), based on mental volition, take a pause. Volitions based on greed, hatred, and delusion stop. ``The relinquishment of all acquisitions (\emph{upadhi}).'' ``Acquisitions'' sometimes refers to the five aggregates affected by clinging, which are entwined with the perception of self as ``me'' or ``mine.'' We cling to the five khandhas with a sense of ownership or identification. Sometimes upadhi are also called the bases or constituents of existence. With disenchantment and dispassion, we relinquish them. There is no need to hold on to them any longer, no need to keep that construct going, no need to cling anymore.

Deep peace also comes with the destruction of craving, ``\emph{taṇhā.}'' Craving is the momentum, the mental push to get something from the sensory world. It is also the momentum to be somebody, to become something, to have an identity, or simply to exist. It is also the craving to annihilate, the aversion to dealing with anything unpleasant. With the destruction of craving, dispassion is experienced. We lose the passion for anything related to the five khandhas, anything related to the conditioned realm. It doesn't mean we cease to exist, but the passion for that which no longer fulfills just fades away. \emph{Nirodha} is translated as cessation or non-arising. Sāriputta said, ``Nibbāna is the cessation of becoming.'' It is the non-arising of the urge to become, to exist. That kind of non-arising is still, peaceful. Dispassion, cessation, Nibbāna: cooling, unbinding, complete freedom, the end of suffering.

This is where contemplation can take us. If we have yet to arrive at that point of deep realization, then bringing these phrases to mind as a possibility, as a perception, is a very powerful contemplation. The teachings of the Buddha and the words that he used to describe the path and its results are powerful, profound, and provide inspiration to keep on the path.
